Indeed, let @@M0@@ be the affine space @@M1@@ whose @@M2@@-points correspond to the sequences @@M3@@, and @@M4@@ the open subset of @@M5@@ corresponding to the @@M6@@-regular sequences. For fixed @@M7@@, the sheaves @@M8@@ corresponding to the @@M9@@-points of @@M10@@ are contained in a flat family over @@M11@@, and by hypothesis @@M12@@ is nonempty. Thus (i) follows from the upper semi-continuity property of the function @@M13@@ (EGA III, 7.7.5). Finally, since every ``sufficiently general'' sequence is @@M14@@-regular and such that @@M15@@ is a subsheaf of @@M16@@ for every @@M17@@, (ii) follows from (i).

\sgaStableTarget{sga6.e13.lem.1-7}\textbf{Lemma 1.7.} \emph{Let @@M18@@ be an exact sequence on @@M19@@. If the Hilbert polynomial of @@M20@@ is written in the form}
@@M21@@
\emph{then}
@@M22@@
Indeed, @@M23@@; therefore @@M24@@, which gives the assertion.

\sgaStableTarget{sga6.e13.prop.1-8}\textbf{Proposition 1.8.} \emph{Let @@M25@@ be a @@M26@@-sheaf on @@M27@@, @@M28@@ and @@M29@@ the Hilbert polynomial of @@M30@@. Then}~:
\begin{enumerate}[label=(\roman*)]
\item \emph{For @@M31@@, @@M32@@.}
\item \emph{@@M33@@ and @@M34@@ is also a @@M35@@-sheaf.}
\end{enumerate}
Indeed, we argue by induction on @@M36@@. When @@M37@@, this is clear since @@M38@@. When @@M39@@, we have an exact sequence
@@M40@@
